\documentclass[11pt]{article}
\usepackage{iftex}

\ifPDFTeX
  \usepackage[T1]{fontenc}
  \usepackage[utf8]{inputenc}
  \IfFileExists{dejavu.sty}{\usepackage{dejavu}\renewcommand{\familydefault}{\rmdefault}}{}
\else
  \usepackage{fontspec}
  \setmainfont{DejaVu Serif}[Scale=0.95]
\fi

\usepackage[a4paper, left=1.4cm, right=1.6cm, top=1.0cm, bottom=1.5cm]{geometry}
\usepackage{enumitem}
\usepackage{calc}
\usepackage{xcolor}
\usepackage[hidelinks]{hyperref}
\pagestyle{empty}
\setlength{\parindent}{0pt}
\setlength{\parskip}{0pt}

\definecolor{linkblue}{HTML}{0563C1}
\newlength{\cvdatewidth}\setlength{\cvdatewidth}{2.05cm}
\newlength{\cvdategap}\setlength{\cvdategap}{0.35cm}
\newlength{\cvlabelwidth}\setlength{\cvlabelwidth}{3.3cm}
\newcommand{\cvnamesize}{\fontsize{26}{30}\selectfont}
\newcommand{\cvsectionsize}{\fontsize{14}{17}\selectfont}
\newcommand{\cvsectionbefore}{14pt}
\newcommand{\cvsectionafter}{8pt}
\newcommand{\cventrygap}{8pt}

\newcommand{\cvlink}[2]{\href{#1}{\textcolor{linkblue}{\underline{#2}}}}
\newcommand{\cvheader}[6]{%
  \begin{center}
    {\cvnamesize\bfseries #1}\\[1pt]
    \cvlink{https://#2}{#2} $|$ \cvlink{https://#3}{#3} $|$ \cvlink{mailto:#4}{#4}\\[4pt]
    #5 $|$ #6
  \end{center}\vspace{-4pt}}

% Highlights details that still need to be filled in
\newcommand{\tbd}[1]{\textcolor{red}{[#1]}}

\newcommand{\cvsection}[1]{%
  \par\vspace{\cvsectionbefore}%
  {\cvsectionsize\bfseries\MakeUppercase{#1}}\par
  \vspace{\cvsectionafter}}

\newlist{cvitems}{itemize}{1}
\setlist[cvitems]{label=-, leftmargin=0.95cm, labelsep=0.5cm,
                  itemsep=1pt, topsep=1pt, parsep=0pt, partopsep=0pt}

\newcommand{\cventry}[4]{%
  \par\noindent
  \begin{minipage}[t]{\cvdatewidth}\raggedright #1\end{minipage}%
  \hspace{\cvdategap}%
  \begin{minipage}[t]{\linewidth-\cvdatewidth-\cvdategap}\raggedright
    \textbf{#2}\if\relax\detokenize{#3}\relax\else\ #3\fi
    \if\relax\detokenize{#4}\relax\else
      \begin{cvitems} #4 \end{cvitems}%
    \fi
  \end{minipage}\par\vspace{\cventrygap}}

\newcommand{\cvskill}[2]{%
  \par\noindent
  \begin{minipage}[t]{\cvlabelwidth}\raggedright\textbf{#1}\end{minipage}%
  \begin{minipage}[t]{\linewidth-\cvlabelwidth}\raggedright #2\end{minipage}%
  \par\vspace{3pt}}


% =====================================================================
\begin{document}

\cvheader{Mathis Burger}
         {mathis-burger.de}
         {github.com/MathisBurger}
         {kontakt@mathis-burger.de}
         {+49 1523 2052626}
         {DE-85055 Ingolstadt}

\cvsection{Education}

\cventry{10/2023 -- 02/2027}{Technische Hochschule Ingolstadt}
  {B.Sc. Business Information Systems (Wirtschaftsinformatik)}
  {\item Current average: 1.2 (German scale, 1.0 = best)
   \item Thesis (expected 01/2027): \textit{Conceptual Design and Prototyping of a Feature-Model-Based Application for Modeling and SAT-Based Verification of Offer Structures} (supervisor: Prof.\ Dr.\ Anatoli Djanatliev).
         Encoding of automotive offer structures as feature models and translation into propositional logic for SAT-based verification.
   \item Student Ambassador}

\cventry{2015 -- 2023}{Werner-Heisenberg-Gymnasium}{Abitur}{}

% ---------------------------------------------------------------------
\cvsection{Awards \& Scholarships}

\cventry{2026 -- present}{Max Weber Programme of the State of Bavaria}{}
  {\item Merit-based scholarship for outstanding students, administered by the
         German Academic Scholarship Foundation (Studienstiftung des deutschen Volkes); Elite Network of Bavaria}

% ---------------------------------------------------------------------
\cvsection{Research Interests}

SAT-based analysis of software dependencies and variability; combining formal reasoning with AI-assisted programming; distributed systems

% ---------------------------------------------------------------------
\cvsection{Projects}

\cventry{10/2026}{SAT-Based Deletion Analysis of Java Code, research prototype}{}
  {\item Encodes Java code dependencies as a feature model and uses SAT solving
        to answer whether code can be deleted and, if not, why
   \item Technical report with specification and worked example: \cvlink{https://github.com/MathisBurger/lbdajcd}{github.com/MathisBurger/lbdajcd}}

\cventry{03/2026 -- 07/2026}{MYnd Learning Platform, university team project}{}
  {\item Platform for individual learning: I was the lead developer in a team of
         13 students, responsible for architecture and backend
   \item Modular monolith following domain-driven design (Java, TypeScript)
   \item Graded 1.0}

\cventry{2024 -- present}{CodeCanvas, open-source learning platform}{}
  {\item Platform for programming assignments with automated, sandboxed code
         execution and individual feedback
   \item Service-based architecture. Deployed at \cvlink{https://code-canvas.app}{code-canvas.app}
   \item Offered to students of my Java tutorials as a practice platform}

% ---------------------------------------------------------------------
\cvsection{Teaching}

\cventry{2026 -- present}{Voluntary Technical Advisor, Technische Hochschule Ingolstadt}{}
  {\item Advisor for a master's project course:
         explaining core distributed-systems concepts and supporting students
         with implementation and debugging}

\cventry{10/2024 -- 07/2025}{Tutor, Technische Hochschule Ingolstadt}{}
  {\item Introduction to Computer Science (winter term 2024/25)
   \item Programming in Java II (summer term 2025), including self-developed
         practice exercises}

% ---------------------------------------------------------------------
\cvsection{Professional Experience}

\cventry{2023 -- present}
  {Corporate Student (Dual Studies), AUDI AG}{Ingolstadt}
  {\item Practical phases in vehicle project management (including software variability management) and digitalisation}

\cventry{01/2026 -- 03/2026}
  {Software Engineer Intern, Volkswagen Group Digital Solutions}{Lisbon}
  {\item Full-stack development (backend and UI) of an application for
         analysing automotive software variability (ELVISNG project)
   \item Part of the Audi dual programme}

\cventry{2021 -- present}
  {Full Stack Software Engineer (Part-time), die NetzWerkstatt}{Remote}
  {\item Data-migration pipelines for legacy systems
   \item Designed and implemented web-widget system using PHP and TypeScript
   \item Developed a web-based tender portal
   \item Paused 09/2023 -- 06/2024 for the dual study programme at AUDI AG}

% ---------------------------------------------------------------------
\cvsection{Other Activities}

\cventry{2026 -- present}{Open-Source Contributor, KERN UX Standard}{}
  {\item Contributed to the KERN React Kit, the official component library implementing the
         KERN UX standard for German public administration}

\cventry{2025 -- present}{Board Member and Officer for Digital Affairs,
  Junge Union Ingolstadt}{}{}

% ---------------------------------------------------------------------
\cvsection{Skills}

\cvskill{Languages}{German (C2), English (C1, IELTS 8.0)}
\cvskill{Programming}{Java, PHP, TypeScript, Kotlin, Python, Go, Rust, SQL}
\cvskill{Methods}{SAT, feature modelling, domain-driven design, Scrum}
\cvskill{Tools}{PostgreSQL, Docker, Git, \LaTeX}

\end{document}
